\section{Introduction}

Plant diseases have historically caused significant ecological and economic damage. In agricultural regions, disease detection often relies on human observation and localized monitoring. In contrast, forested and remote regions lack continuous surveillance. Advances in satellite imagery and environmental data collection enable long-term observation of vegetation behavior across large geographic regions, but analyzing such data effectively requires approaches that go beyond traditional single-machine scripts and instead follow scalable data processing paradigms~\cite{verbesselt2010, xie2025}.

The Mountain Pine Beetle (MPB), \textit{Dendroctonus ponderosae}, provides a compelling case study. Over the past two decades, MPB outbreaks have caused the mortality of hundreds of millions of trees across western North America, profoundly disrupting forest carbon storage, water cycling, and wildfire dynamics. Research has shown that vegetation spectral signatures, particularly NDVI and moisture-sensitive indices, begin to diverge from healthy seasonal patterns weeks to months before visible symptoms appear~\cite{jamali2024, lasaponara2024}. Detecting these subtle shifts at scale requires a Big Data infrastructure capable of ingesting, indexing, and querying large volumes of multi-temporal raster data.

The primary objective of this project is to analyze historical vegetation trends in regions where plant diseases have been reported. For selected disease events, vegetation behavior is studied over multiple years before the documented detection date. The analysis seeks to identify deviations or anomalies in vegetation patterns that may precede disease outbreaks. The project follows a data-driven and exploratory approach and does not assume the existence of strong predictive signals. Rather than serving directly as an operational warning system, the framework evaluates whether anomalous vegetation patterns were present in the period leading to documented outbreaks, acknowledging the possibility of false positives while aiming to support earlier intervention.

This project contributes: (1) a fully operational, cloud-native ETL pipeline for COG raster ingestion into AWS S3 with PostGIS-indexed spatiotemporal metadata; (2) a PySpark distributed processing stack for NDVI/NDMI computation and monthly temporal compositing at scale; (3) a Z-score anomaly detection system benchmarked against a multi-year historical baseline; (4) three independent scalability demonstrations across datasets of differing resolution, volume, and temporal extent; and (5) output products including raster anomaly maps, Parquet event datasets, and animated spectral timelapses. The project is structured so that it can be extended to significantly larger datasets without fundamental changes to the pipeline. Evaluation demonstrates that GIST-indexed queries reduce latency by 15--25$\times$ over sequential scans, PySpark throughput scales approximately linearly with data volume, and the Z-score detector achieves approximately 71\% recall in identifying vegetation stress events corresponding to documented MPB activity.

The remainder of this report is organized as follows. Section~2 reviews related work grouped into four categories. Section~3 describes the datasets and study contexts. Section~4 details data management and indexing. Section~5 covers the distributed transformation pipeline. Section~6 presents anomaly detection methodology. Section~7 reports the evaluation. Section~8 describes the visualization of outputs. Section~9 concludes the paper and discusses future directions.